\section*{METAMATHEMATICS}
\phantomsection
\addcontentsline{toc}{section}{METAMATHEMATICS}

The absence of contradiction has at all times been considered as a sine qua non condition for all mathematics, and from the time of Aristotle, logic was sufficiently developed so that it was perfectly well known that from an inconsistent theory anything could be deduced. The proofs of ``existence'', considered as indispensable since Antiquity, visibly had no other aim than to guarantee that the introduction of a new concept did not risk entailing a contradiction, particularly when this concept was too complicated to fall immediately under ``intuition''. We have seen how this requirement became more pressing with the arrival of the axiomatic point of view in the XIXth century and how the construction of arithmetical ``models'' responded to it. But arithmetic itself, could it be inconsistent? A question that without doubt would not have been thought of before the end of the XIXth century, so much did the whole numbers appear to belong to what was most certain in our intuition; but after the ``paradoxes'' everything seemed to have been questioned, and one can understand that the feeling of insecurity that they created had brought mathematicians, around 1900, to bear with more care on the problem of the non-contradiction of arithmetic, so as to save at least classical mathematics from shipwreck. So this problem is the second of those that Hilbert enumerated in his famous lecture at the international Congress of 1900 ({[}163 a{]}, v. III, pp.~229-301). In doing this, he laid down a new principle that was to have great repercussions: whereas in traditional logic the non-contradiction of a concept only made it ``possible'', it was equivalent for Hilbert (at least for mathematical concepts defined axiomatically) to the existence of the concept. This implied apparently the necessity of proving a priori the non-contradiction of a mathematical theory before even being able to develop it legitimately; it is just in this way that H. Poincaré understands it when, to make a breach in formalism, he takes to his account the idea of Hilbert, underlining with a mischievous pleasure the extent to which the formalists were far, at that stage, from being able to achieve it ({[}251 c{]}, p.~163). We will see later how Hilbert was to pick up the gauntlet; but before that, it must be noted here that, under the influence of his authority and that of Poincaré, the requirements set down by this latter were to be for a long time accepted without reservation as well by the formalists as by their adversaries. One consequence was the belief, very widespread among the formalists, that the theory of proof of Hilbert was an integral part of mathematics, of which it constituted the indispensable prolegomena. This dogma does not seem justified to us, \(^{67}\) and we consider that the intervention of metamathematics in the exposition of logic and mathematics can and should be reduced to the very elementary part that is concerned with the handling of abbreviating symbols and criteria of deduction. It is thus not a case, contrary to what Poincaré thinks, of ``claiming the freedom from contradiction'', but rather to consider, with Hadamard, that the absence of contradiction, even when it is not proven, is noted ({[}12{]}, p.~270).

It remains for us to give a brief historical sketch of the efforts of Hilbert and his school, and to retrace summarily, not only the evolution that was finally to lead to the negative result of Gödel and justify a posteriori the scepticism of Hadamard, but also all the progress that followed on from this, concerning the knowledge of the mechanism of mathematical reasoning, and which makes modern metamathematics an autonomous science of incontestable interest.

Already in 1904, in a lecture at the international Congress ({[}163 c{]}, pp.~247-261), Hilbert attacks the problem of consistency in arithmetic. He observes first that there is no question of proving it by making use of a model, \(^{68}\) and outlines the principles of another method: he proposes considering the true propositions of formalised arithmetic as collections of signs without meaning, and to prove that in using the rules governing the formation and linkage of these collections, it is impossible to obtain a collection which is a true proposition and whose negation is also a true proposition. He sketches even a proof of this nature for a formalism less extensive than that of arithmetic; but, as was observed by H. Poincaré soon after ({[}251 c{]}, p.~185), this proof uses essentially the principle of recurrence, and so appears to rely on a vicious circle. Hilbert did not answer this criticism immediately, and about fifteen years passed by without anyone attempting to develop his ideas; it is in 1917 only that (motivated by the wish to answer the attacks of the intuitionists) he harnesses himself again to the problem of the foundations of mathematics, which he will not cease from now on to consider until the end of his scientific career. In his work on this subject, which is spread out from 1920 to 1930 approximately, and in which a whole school of young mathematicians take part (Ackermann, Bernays, Herbrand, von Neumann), Hilbert brings out little by little the principles of his ``theory of proof'' in a more precise way: acknowledging implicitly the validity of the criticism of Poincaré, he admits that in metamathematics, the arithmetical reasoning that is used can only be based on our intuition of the whole numbers (and not on formalised arithmetic); to do this it appears to be essential to restrict this reasoning to ``finite procedures'' (``finite Prozesse'') of a type allowed by the intuitionists: for example, a proof by contradiction cannot prove the metamathematical existence of a collection or sequence of collections, there must be an explicit law of construction. \(^{69}\) On the other hand Hilbert enlarges in two directions his initial programme: not only does he tackle the consistency of arithmetic, but he aspires also to prove the consistency of the theory of real numbers and even that of the theory of sets; \(^{70}\) further, to the problems of consistency come to be added those of independence of axioms, of completeness and of decision. We will pass quickly in review these various questions and point out the principal research to which they have given rise.

The proof of the independence of a system of propositions \(A_{1}, A_{2}, \ldots, A_{n}\) consists in showing that, for each index i, \(A_{i}\) is not a theorem in the theory \(T_{i}\) obtained by taking as axioms the \(A_{j}\) with index \(j \neq i\) . This point will be established if one knows a consistent theory \(T_{i}^{\prime}\) in which the \(A_{j}(j \neq i)\) are theorems, as well as ``not \(A_{i}\) ''; thus one can consider the problem in two ways depending on whether or not one allows that certain theories (such as arithmetic or the theory of sets) are consistent. In the second case, one has to deal with a problem of ``absolute'' consistency. On the contrary, the first type of problem reduces to a problem of ``relative'' consistency, by means of a construction of appropriate ``models'', and many proofs of this kind have been thought up, well before mathematics had taken up a completely formalised aspect: it is enough to remember the models of non-Euclidean geometry (see p.~133), the questions of the independence of the axioms of elementary geometry dealt with by Hilbert in the ``Grundlagen der Geometrie'' {[}163 c{]} as well as the work of Steinitz on the axiomatisation of Algebra and those of Hausdorff and his successors on that of Topology.

A theory T is said to be complete if, for every proposition A of T (not containing any letter other that the constants of T), one of the propositions A, (not A), is a theorem of \(T.^{71}\) . If one puts on one side certain very rudimentary formalisms, of which the completeness is easily proved ({[}164{]}, p.~35), the results obtained in this area are essentially negative; the first in time, is due to K. Gödel {[}130 a{]} who showed that, if T is consistent and if the axioms of formalised arithmetic are theorems of T, then T is not complete. The fundamental idea of his ingenious method consists in establishing a bijective correspondence (of course, by means of ``finite processes'') between the metamathematical statements and certain propositions in formalised arithmetic; we will restrict ourselves to a broad sketch. \(^{72}\) For each collection A which is a term or a relation of T, one begins by associating (by a procedure which is an explicit construction, which can be applied quasi mechanically) a whole number \(g(A)\) , in a bijective way. In the same way, to each proof D of T (considered as a succession of collections) one can associate a whole number \(h(D)\) in a bijective way. Finally, one can give a procedure for constructing explicitly a relation \(P(x,y,z)\) of T, \(^{73}\) such that, in T, \(P(x,y,z)\) requires that x,y,z are whole numbers, and satisfies the following two conditions:

\(1^{\circ}\) if \(D\) is a proof of \(A(\lambda)\) , where \(A(x)\) is a relation of \(\mathcal{T}\) , and \(\lambda\) is an explicit whole number (that is to say a term of \(\mathcal{T}\) which is a whole number), then \(P(\lambda, g(A(x)), h(D))\) is a theorem of \(\mathcal{T}\) ;

\(2^{\circ}\) if the explicit whole number \(\mu\) is not of the form \(h(D)\) , or if \(\mu = h(D)\) and if \(D\) is not a proof of \(A(\lambda)\) , then (not \(P(\lambda, g(A(x)), \mu)\) ) is a theorem of \(\mathcal{T}\) .

Then let \(S(x)\) be the relation (not \((\exists z)P(x,y,z)\) ), and let \(\gamma = g(S(x))\) , which is a term of \(\mathcal{T}\) . If \(\mathcal{T}\) is consistent, there is no proof of the proposition \(S(\gamma)\) in \(\mathcal{T}\) . In fact, if \(D\) were such a proof, \(P(\gamma,g(S(x)),h(D))\) would be a theorem of \(\mathcal{T}\) ; but this relation is none other than \(P(\gamma,\gamma,h(D))\) , and it follows that \((\exists z)(P(\gamma,\gamma,z)\) would also be a theorem of \(\mathcal{T}\) ; as this last relation is equivalent to (not \(S(\gamma)\) ), \(\mathcal{T}\) would be inconsistent. On the other hand what has just been said shows that for every explicit whole number \(\mu\) , (not \(P(\gamma,\gamma,\mu)\) ) is a theorem of \(\mathcal{T}\) . The result is that there is no proof, in \(\mathcal{T}\) , of (not \(S(\gamma)\) ), for this last relation is equivalent to \((\exists z)P(\gamma,\gamma,z)\) and the existence of a whole number \(\mu\) such that \(P(\gamma,\gamma,\mu)\) would require that \(\mathcal{T}\) were inconsistent, by virtue of what comes before. \(^{74}\) This metamathematical theorem of Gödel has been subsequently generalised in various directions ({[}181{]}, chap.~XI). \(^{75}\)

The relation \(S(\gamma)\) of T of which it is shown thus that there is no proof in T either of this relation, or of its negation, is obviously made up for the needs of the cause and does not tie up in any natural way with any mathematical problem. Much more interesting is the fact that if \(\mathcal{T}\) denotes the theory of sets (with the system of axioms of von Neumann-Bernays), neither the continuum hypothesis nor its negation are provable in \(\mathcal{T}\) . This remarkable result was established in two stages: in 1940, Gödel proved that the theory obtained by adjoining to \(\mathcal{T}\) the continuum hypothesis \(2^{\aleph_0} = \aleph_1\) was not inconsistent {[}130 b{]}; then, in 1963, P. Cohen proved the same when one adjoins to \(\mathcal{T}\) the relation \(2^{\aleph_0} = \aleph_2\) (or \(2^{\aleph_0} = \aleph_n\) for an arbitrary whole number \(n > 1\) ) {[}67{]}.

The decision problem (``Entscheidungsproblem'') is without doubt the most ambitious of all those which are stated in metamathematics: it consists in knowing whether, for a given formal language, one can imagine a ``universal procedure'' which is quasi mechanical which, when applied to any relation of the formalism in question, indicates in a finite number of operations whether that relation is true or not. \(^{76}\) The solution of this problem was already part of the grand design of Leibniz, and it appears that at one time the school of Hilbert believed its realisation was quite close. It is a fact that one can describe such procedures for formalisms containing few primitive symbols and axioms ({[}181{]}, pp.~136-141). But the efforts made to make precise the decision problem, by outlining exactly what is meant by ``universal procedure'' have until now resulted only in negative results ({[}181{]}, pp.432-439). Besides, the solution of the decision problem for a theory T entrains immediately the knowledge of whether T is consistent or not, since it suffices to apply the ``universal procedure'' to a relation of T and its negation; and we will see that it is impossible to solve the question in this way for the usual mathematical theories. \(^{77}\)

It is indeed in the question of the consistency of mathematical theories --- the origin and very heart of Metamathematics --- that the results have revealed themselves to be the most deceiving. During the years 1920-1930, Hilbert and his school had developed new methods for attacking these problems; after having proved the consistency of partial formalisms, covering part of arithmetic (cf.~{[}158{]}, {[}324 d{]}), they believed they had reached the goal and proved, not only the consistency of arithmetic, but also that of the theory of sets, when Gödel, relying on the incompleteness of arithmetic, deduced from it the impossibility of proving, by the ``finite procedures'' of Hilbert, the consistency of every theory T containing this latter. \(^{78}\)

The theorem of Gödel does not however shut the door completely on attempts to prove consistency, so long as one abandons (at least partially) the restrictions of Hilbert concerning ``finite procedures''. It is thus that Gentzen in 1936 {[}127{]}, succeeds in proving the consistency of formalised arithmetic, by using ``intuitively'' transfinite induction up to the countable ordinal \(\varepsilon_{0}\) .⁷⁹ The value of the ``certainty'' that one can attach to such reasoning is without doubt less convincing than for that which satisfies the initial requirements of Hilbert, and is essentially a matter of personal psychology for each mathematician; it remains no less true that similar ``proofs'' using ``intuitive'' transfinite induction up to a given ordinal, would be considered important progress if they could be applied, for example, to the theory of real numbers or to a substantial part of the theory of sets.

